% VI38-DETAILED-PROBLEM-03
\begin{problem}[VI.38.P3: Quotient maps and closed immersion]\label{prob:vi38-03}
Let \(S\twoheadrightarrow T\) be a surjective map of standard graded \(k\)-algebras.  Prove the induced morphism \(\Proj T\to\Proj S\) is a closed immersion without appealing to the global quotient theorem of VI/37.
\end{problem}

\ifdefined\FullProblemDossiers
\begin{solution}
Write \(T=S/I\) for the homogeneous kernel \(I\).  Because the map is surjective and
both rings are standard graded,
\[
T_+=\varphi(S_+)T.
\]
No point \(\mathfrak q\in\Proj T\) contains \(T_+\), so the Proj morphism is defined
everywhere.

Let \(f\in S_+\) be homogeneous and write \(\bar f\) for its image in \(T\).  On the
standard chart \(D_+(\bar f)\), localization of the quotient gives
\[
T_{\bar f}\cong S_f/IS_f.
\]
Taking degree-zero parts yields
\[
T_{(\bar f)}
\cong
S_{(f)}/(IS_f)_0.
\]
Thus
\[
S_{(f)}\twoheadrightarrow T_{(\bar f)}
\]
is surjective.  Contravariantly,
\[
D_+(\bar f)=\Spec T_{(\bar f)}
\longrightarrow
D_+(f)=\Spec S_{(f)}
\]
is an affine closed immersion.

These maps are restrictions of the global Proj morphism and agree on the standard
overlaps by localization in stages. Since the \(D_+(f)\) form an affine open cover
of \(\Proj S\), the affine-local criterion for closed immersions gives
\[
\boxed{\Proj T\hookrightarrow\Proj S}.
\]
Moreover, on each chart its image is cut out by \((IS_f)_0\).  Hence the global
image is the closed subscheme determined by the homogeneous kernel \(I\), recovering
the quotient description chartwise without invoking VI/37's global theorem.
\end{solution}
\fi
